% ===========================================================================
\section{Verificação e reprodutibilidade}\label{sec:verificacao}
% ===========================================================================

As figuras deste artigo são calculadas a partir das definições por um programa
curto em Python, e as tabelas por outro; ambos acompanham o texto. A aritmética da
Seção~\ref{sec:blocos} e as cotas da Seção~\ref{sec:geometria} foram também
verificadas de forma mecânica em Lean~4. Esta seção descreve o que foi calculado,
como, e o que a verificação formal cobre.

\subsection{Os programas}

\begin{sloppypar}
O material suplementar, disponível em
\url{https://github.com/tiagobandeira/analise-espectral-primos}, contém:
\end{sloppypar}

\begin{center}
\footnotesize
\begin{tabular}{llr}
\toprule
arquivo & o que calcula & tempo \\
\midrule
\texttt{figuras\_artigo.py} & as seis figuras (PDF vetorial) & $\approx 8$ min \\
\texttt{verificar\_teoremas.py} & as tabelas e as verificações V1--V5 & $\approx 2$ min \\
\bottomrule
\end{tabular}
\end{center}

O primeiro programa não depende de nenhuma biblioteca além de \texttt{numpy} e
\texttt{matplotlib}, e recalcula cada painel a partir da definição do sinal
\eqref{eq:Sm}: os logaritmos de módulo são somados termo a termo, as transformadas
são calculadas com \texttt{rfft} e as frequências convertidas por
$x=\operatorname{round}(e^{2\pi f})$. O segundo verifica por força bruta, para
blocos até $k=22$, que o menor fator primo de todo composto de $A_k$ está em
$\Seed_k$ (Lema~\ref{lema:bloco}), que o par de vizinhos inteiros de todo primo do
bloco é $\Seed_k$-suave (Lema~\ref{lema:qmaisum}) e que o mínimo de $\rB$ é
atingido em $q_k$ (Corolário~\ref{cor:min}).

\subsection{Os números das tabelas}

As Tabelas~\ref{tab:crivo}, \ref{tab:rho} e \ref{tab:lei} contêm os valores
medidos. A Tabela~\ref{tab:rho} é inteiramente explicada pelo
Teorema~\ref{teo:exato}: cada entrada é $g(q_k)$, calculada em precisão dupla, e a
coluna da cota é a expressão do Teorema~\ref{teo:cota}. A Tabela~\ref{tab:lei}
contém o par gargalo de cada semente, obtido por busca direta sobre os pares de
primos consecutivos.

A Tabela~\ref{tab:janela} e a Tabela~\ref{tab:gemeos} são medições, não
identidades: nelas há erro de leitura, e é justamente esse erro que a
Proposição~\ref{prop:janela} explica. Registramos nas duas tabelas o valor medido,
sem suavização.

\subsection{A verificação formal}

O núcleo da Seção~\ref{sec:geometria} e a aritmética da Seção~\ref{sec:blocos}
foram formalizados em Lean~4 sobre a biblioteca Mathlib \cite{mathlib}, sem
\texttt{sorry} e sem axiomas além dos da própria biblioteca. O Apêndice~\ref{ap:lean}
traz os enunciados exatos.

\begin{table}[htbp]
\centering
\small
\begin{tabular}{lll}
\toprule
resultado & declaração em Lean & status \\
\midrule
Lema~\ref{lema:bloco} (composto tem fator pequeno) &
\texttt{existe\_fator\_primo\_na\_semente} & provado \\
Teorema~\ref{teo:crivo} (crivo por blocos) &
\texttt{bloco\_primo\_iff} & provado \\
Teorema~\ref{teo:irred} (irredutibilidade) &
\texttt{logRedutivel\_iff} & provado \\
Lema~\ref{lema:vizinho} (vizinho inteiro) &
\texttt{log\_dist\_ge} & provado \\
Teorema~\ref{teo:cota} (cota uniforme) &
\texttt{cota\_rho} & provado \\
monotonia da cota em $k$ (auxiliar) &
\texttt{cota\_antitone} & provado \\
$\log(1+1/q)\ge 1/(q+1)$ &
\texttt{inv\_succ\_le\_log\_ratio} & provado \\
igualdade em $q+1$ &
\texttt{log\_dist\_succ\_eq} & provado \\
\bottomrule
\end{tabular}
\caption{O que a formalização cobre. As declarações residem em
\texttt{Espectral/Bloco.lean} e \texttt{Espectral/Rho.lean}, e compilam contra
\texttt{leanprover/lean4:v4.37.0-rc1} com Mathlib.}
\label{tab:lean}
\end{table}

O que a formalização \emph{não} cobre, e por quê:

\begin{itemize}[leftmargin=1.4em,itemsep=2pt]
\item A igualdade do Teorema~\ref{teo:exato} na forma completa exigiria definir
$\rB$ como mínimo sobre um conjunto de suaves; formalizamos as duas metades que a
demonstração usa — a cota inferior para todo inteiro (Lema~\ref{lema:vizinho}) e o
fato de a cota ser atingida em $q+1$ —, que é o conteúdo matemático. O passo
final, ``logo o mínimo é esse valor'', é imediato no papel.
\item A suavidade de $q+1$ (Lema~\ref{lema:qmaisum}) é aritmética elementar mas
exige a maquinaria de fatoração de Mathlib; deixamo-la no papel.
\item As Proposições~\ref{prop:gemeo} e \ref{prop:janela} são afirmações sobre
integrais de somas finitas de cossenos. Formalizá-las exigiria a teoria de
integrais de Riemann para funções oscilatórias, com custo desproporcional ao
conteúdo, que é uma desigualdade triangular e $|F(\alpha)|\le2/|\alpha|$.
\end{itemize}

\subsection{Como reproduzir}

\begin{center}
\begin{minipage}{0.95\textwidth}
\footnotesize
\begin{verbatim}
# figuras (Python 3.11+, numpy, matplotlib)
python artigo/experimentos/figuras_artigo.py

# tabelas e verificações por força bruta
python artigo/experimentos/verificar_teoremas.py

# formalização (Lean 4 + Mathlib)
cd formalizacao && lake exe cache get && lake build
\end{verbatim}
\end{minipage}
\end{center}

Nenhum dos cálculos usa amostragem aleatória: todos são determinísticos, e os
valores das tabelas são reprodutíveis bit a bit em precisão dupla. Os arquivos
\texttt{.json} produzidos pelos programas contêm os valores brutos de cada medida
citada no texto.
